%قبل از شروع مسئله ترجیحاً چیزی ننویسید. اما اگر خواستید بنویسید، در قالب لیستی (فقرات و شمارش) بنویسید.
%برای اضافه کردن نمره، بعدی از شروع‌مسئله، [x] را اضافه کنید و به جای x نمره مورد نظر را قرار دهید. برای زیرمسئله‌ها هم به همین صورت است.
%\شروع‌مسئله[10]
\شروع‌مسئله [15]
%\متن‌سیاه{سوال عملی سوم:}\\
%\متن‌سیاه{دراپ‌تست}\\

%\متن‌سیاه{ورودی:}\\
%در خط اول ورودی، ابتدا k و سپس m داده می‌شود.\\
%\متن‌سیاه{خروجی:}\\
%در تنها سطر خروجی، تعداد آیفون‌های مورد نیاز می‌آید. اگر در هیچ صورتی چنین کاری ممکن نبود ۱- چاپ می‌شود.\\
%\متن‌سیاه{محدودیت‌ها:}\\
%$1\leq k\leq 10^4$

%$m\leq10^9$
%\\
%\textbf{\متن‌سیاه{ورودی و خروجی نمونه}\\
%% تست اول
%\begin{latin}
% \begin{tabularx}{\textwidth}{|X | X|}
% \hline
% \textbf{stdin} &
% \textbf{stdout}\\
% \hline
% 100 100&
% 1\\
% \hline
% 
% \end{tabularx}
%\end{latin}
%
%% تست دوم
%
%\begin{latin}
% \begin{tabularx}{\textwidth}{|X | X|}
% \hline
% \textbf{stdin} &
% \textbf{stdout}\\
% \hline
% 340 20&
% 3\\
% \hline
% 
% \end{tabularx}
%\end{latin}}
% صورت اصلی مسئله را در این‌جا بنویسید.
فرانک لمپارد، سرمربی محبوب تیم فوتبال چلسی، قصد دارد تا تمرینات فشردهای را برای بازیکنان این تیم، تدارک ببیند.
زمین تمرین تیم چلسی، به شکل یک آرایه از اعداد است! لمپارد میخواهد قسمتی از این زمین را که به صورت یک پیشوند از
این آرایه است (یعنی تعدادی خانهی متوالی، با شروع از خانهی اول – این تعداد میتواند صفر نیز باشد) را برای تمرین
کارهای دفاعی مدافعان (آزپیلیکوئتا، زوما، رودیگر و مارکوس آلونسو) تخصیص دهد. و همچنین میخواهد قسمتی از این
زمین را که به صورت یک پسوند از این آرایه است (یعنی تعدادی خانهی متوالی، با شروع از خانهی آخر آرایه – این تعداد میتواند صفر نیز باشد) را برای تمرین کارهای هجومیِ مهاجمانش (ویلیان، پدرو، پولیسیچ و تامی آبراهام) اختصاص دهد .
واضح است که این دو قسمت از زمین تمرین، نباید با هم اشتراک داشته باشند.

متاسفانه، زمین تمرین از ناهمواریهایی تشکیل شده است. میزان مناسب بودن یک قسمت از زمین تمرین را برابر با \متن‌لاتین{xor}
تمام اعداد موجود در آن قسمت از زمین، تعریف میکنیم. و میزان اثربخشی تمرین را برابر با \متن‌لاتین{xor} میزان مناسب بودن دو قسمت از زمین، تعریف میکنیم.

حال لمپارد میخواهد بیشترین میزان اثربخشی ممکن که میتواند برای تمرینش بهدستبیاورد را پیدا کند. این کار را انجام دهید.

\متن‌سیاه{ورودی:}\\
در خط اول، $n$ تعداد اعداد آرایهی زمین تمرین – داده میشود.
سپس در خط بعد، $n$ عدد $a_1, a_2, , ... ,a_n$ ورودی داده می شود، که به ترتیب از چپ به راست، اعداد آرایه ی زمین تمرین
هستند.

\متن‌سیاه{خروجی:}\\
در تنها خط خروجی، بیشترین میزان اثربخشی ممکن را چاپ کنید.

\متن‌سیاه{محدودیت‌ها:}\\
$1 \leq n \leq 100\ 000$ \\
$1 \leq a_i \leq 10^{12}$



%$m\leq10^9$

%\textbf{\متن‌سیاه{ورودی و خروجی نمونه}\\

\متن‌سیاه{ورودی و خروجی نمونه:}\\


%% تست اول
\begin{latin}
	\begin{tabularx}{\textwidth}{|X | X|}
		\hline
		\textbf{stdin} &
		\textbf{stdout}\\
		\hline
		3
		
		3 2 1
		
		&
		1
		\\
		\hline
		
	\end{tabularx}
\end{latin}

\solution{
	
	به کمک \متن‌لاتین{fenwick tree} قابل حل است.
	
}

%%%%%%%%%%%%%%%%%%%%%%%%%%%%%%%%%%%%%%%%%%%%%%%%%%%%%
% برخی دستورات مفید برای تمرین‌ها. %
%%%%%%%%%%%%%%%%%%%%%%%%%%%%%%%%%%%%%%%%%%%%%%%%%%%%%

% فرمول‌های نمونه:
%$f(x) = \Theta(x)$
%$T(x, y) \leq \log(x) + 2^y + \binom{2n}{n}$
% {\tt code\_function}


% لیست موارد بدون شماره‌گذاری:
%\شروع{فقرات}
% \فقره مورد اول 
% \شروع{فقرات}
% \فقره مورد اول داخلی 
% \شروع{فقرات}
% \فقره اولین مورد خیلی داخلی 
% \فقره دومین مورد خیلی داخلی 
% \پایان{فقرات}
% \فقره مورد دوم داخلی 
% \پایان{فقرات}
% \فقره مورد دوم 
%\پایان{فقرات}

% لیست موارد شماره‌گذاری شده:
%\شروع{شمارش}
% \فقره یک مورد 
% \فقره یک مورد دیگر 
% \فقره باز هم یک مورد دیگر
%\پایان{شمارش}

% برای نوشتن یک فرمول چند خطی که یک شماره بخوره. برای حالت‌های دیگر به سایت زیر مراجعه کنید:
% https://www.overleaf.com/learn/latex/Aligning_equations_with_amsmath
%\begin{align} 
% \begin{split}
% (x+y)^3 &= (x+y)^2(x+y) \\
% &= (x^2+2xy+y^2)(x+y) \\
% &= (x^3+2x^2y+xy^2) + (x^2y+2xy^2+y^3) \\
% &= x^3+3x^2y+3xy^2+y^3
% \end{split} 
%\end{align}

% نوشتن ماتریس و آرایه دوبعدی
%\begin{align}
% A = 
% \begin{bmatrix}
% A_{11} & A_{21} \\
% A_{21} & A_{22}
% \end{bmatrix}
%\end{align}

% اضافه کردن تصویر با شرح. برای اضافه کردن تصویر بدون شرح، فقط دستور درج‌تصویر را استفاده کنید.
%\شروع{شکل}[h!]
%\تنظیم‌ازوسط
% \درج‌تصویر[پهنا=0.5\پهنای‌متن]{img.jpg}
%\شرح{شرح تصویر}
%\برچسب{تصویر:۱}
%\پایان{شکل}
%\پایان{زیرمسائل}

%کد زیر را ببینید.
%%\lstinputlisting{source_filename.py}
%\begin{latin}
%\begin{lstlisting}
%# Import the modules
%import sys
%import random
%
%ans = True
%
%while ans:
%question = raw_input("Ask the magic 8 ball a question: (press enter to quit) ")
%
%answers = random.randint(1,8)
%
%if question == "":
%sys.exit()
%
%elif answers == 1:
%print "It is certain"
%
%elif answers == 2:
%print "Outlook good"
%
%elif answers == 3:
%print "You may rely on it"
%
%elif answers == 4:
%print "Ask again later"
%
%elif answers == 5:
%print "Concentrate and ask again"
%
%elif answers == 6:
%print "Reply hazy, try again"
%
%elif answers == 7:
%print "My reply is no"
%
%elif answers == 8:
%print "My sources say no"
%\end{lstlisting}
%\end{latin}